\section{Anthropic, Neo Labs y la diferenciación entre empresas de modelos}

La sección anterior desglosó OpenAI en productos y flujo de efectivo; esta sección examina la diferenciación entre empresas de modelos. La diferencia entre Anthropic y OpenAI no se reduce a que sus modelos tengan nombres distintos: difieren la estructura de clientes, la línea de productos, el control del margen bruto y los supuestos de sus proyecciones financieras. Neo Labs representa otro tipo de diferenciación: equipos de investigación que salen de las grandes empresas de modelos y siguen apostando por Continuous Learning, el entrenamiento personalizado y el ecosistema de código abierto.

\begin{figure}[H]
\centering
\includegraphics[width=0.94\textwidth]{openai-vs-anthropic.png}
\caption{OpenAI vs Anthropic: una empresa de productos 2C frente a una empresa 2B/Coding. Diagrama conceptual propio, elaborado a partir de la conversación de 00:28:25--00:31:25.}
\end{figure}

\begin{knowledgebox}{Lectura de la figura: las dos empresas no son copias del mismo negocio}
La ventaja de OpenAI se inclina hacia el posicionamiento entre los consumidores, el asistente personal y el punto de entrada empresarial; más de la mitad de los ingresos de Anthropic proviene del segmento 2B, y la empresa pone más énfasis en Coding, Finance, los modelos de Agent y los trabajadores del conocimiento en las empresas. Ambas pueden ir muy bien, pero difieren en la calidad de sus ingresos, sus vías de venta y la forma en que organizan sus productos.
\end{knowledgebox}

\subsection{El crecimiento de Anthropic y el control del margen bruto}

Freda menciona que los ingresos anualizados de Anthropic crecieron muy rápido, de 1,000 millones a 7,000 millones, y que en los últimos meses su incremento mensual de ingresos anualizados se ha acercado en magnitud al de OpenAI. Más importante es la mejora del margen bruto: mediante la arquitectura del modelo, el auto scaling, la optimización de la inferencia y otros medios, el margen bruto de una empresa de modelos no está condenado a ser bajo. Ella considera que no es inimaginable que, a largo plazo, el margen bruto de las empresas de modelos llegue a 70\%--80\%.

Aquí hay que atender a la diferencia en los supuestos de las proyecciones financieras. Según ella, la proyección de Anthropic parece suponer que después de 2028 el cómputo total ya no crecerá de forma considerable, de modo que la desaceleración del crecimiento de los costos de entrenamiento eleva el margen de utilidad; la proyección de OpenAI, en cambio, parece suponer que el ROI de los modelos aumenta año con año, es decir, que cada peso invertido en entrenamiento podrá recuperarse en un múltiplo mayor en el futuro. Ninguno de los dos supuestos es un hecho en sí mismo: son modelos sobre los que conviene que los inversionistas pregunten a fondo.

\begin{importantbox}{Las dos palancas de las proyecciones financieras de una empresa de modelos}
La primera es la palanca de costos: si el crecimiento de los costos de entrenamiento e inferencia se desacelera. La segunda es la palanca de ingresos: si cada generación de modelos aporta un ROI mayor. Cada empresa puede construir su relato de mejora del margen de utilidad con una palanca distinta.
\end{importantbox}

\subsection{Neo Labs: la lógica de financiamiento de la siguiente generación de empresas de modelos}

Esta subsección amplía la perspectiva de los dos líderes, OpenAI y Anthropic, a los nuevos laboratorios. Si, como se dijo antes, las empresas de modelos se diferenciarán, la pregunta de Neo Labs es: cuando las empresas de modelos punteras ya son muy fuertes, ¿en qué pueden apoyarse los nuevos equipos para obtener financiamiento y diferenciarse?

En el programa se mencionan nombres conocidos como SSI y Thinking Machines, y también que en los últimos meses varias Neo Labs más han obtenido financiamiento; muchas fueron fundadas por personas con trayectoria académica y los montos recaudados son muy grandes. La duda de Freda es: ¿los modelos grandes terminarán por concentrarse en una o dos empresas punteras, mientras las demás se vuelcan al código abierto? El primer producto de Thinking Machines, Thinker, se describe como una herramienta que ayuda a los usuarios a hacer entrenamiento personalizado de modelos de código abierto; en el fondo, es una apuesta a que el mercado de código abierto será cada vez más grande.

\begin{knowledgebox}{Términos clave: Neo Labs y personalización de código abierto}
\begin{tabularx}{\textwidth}{p{0.24\textwidth}p{0.32\textwidth}X}
\toprule
Término & Explicación breve & Relación con este episodio \\
\midrule
Neo Labs & Nueva generación de laboratorios de modelos o de investigación, a menudo fundados por talento proveniente de grandes empresas o de la academia & Aprovechan las oportunidades del siguiente paradigma, la personalización de código abierto y los modelos especializados. \\
Continuous Learning & Aprendizaje continuo: mejorar el modelo a partir de nuevas interacciones y datos & En el programa se menciona que SSI podría estar explorando esta dirección. \\
Entrenamiento personalizado de código abierto & Entrenamiento específico para el usuario sobre la base de un modelo de código abierto & Thinker, de Thinking Machines, se usa para ilustrar esta dirección. \\
Valuación con OpenAI como referencia & Usar la valuación total de OpenAI como ancla relativa para otras empresas de modelos & Referencia psicológica habitual en el financiamiento de empresas de modelos en el mercado primario. \\
\bottomrule
\end{tabularx}
\end{knowledgebox}

\subsection{Resumen de la sección}

Anthropic y Neo Labs muestran que las empresas de modelos ya se han diferenciado: unas son el punto de entrada para los consumidores, otras se enfocan en empresas y Coding, otras en la personalización de código abierto y otras apuestan por el siguiente paradigma. El mercado de capitales seguirá usando a OpenAI como ancla, pero las diferencias reales provienen de la estructura de clientes, el control del margen bruto y la línea de productos.
